\documentclass[10pt,a4paper]{amsart}
\usepackage{amsmath,amssymb,amsthm,graphicx}
% Retained preamble for the Stacks Project chapter snapshot in this directory.
% The chapter's own preamble.tex is a copy of the shared Stacks Project preamble; the block
% below is sliced MECHANICALLY, line for line, from it. Nothing was added, removed or
% redefined; the wrapper class, the xr-hyper/hyperref wiring and the externaldocument list
% are omitted because the assembler supplies the document class.
%
% lines taken from preamble.tex: 3, 5, 9, 11, 12, 156, 157, 158, ... (51 lines total)
%
% package coverage restored from the chapter's own preamble (2026-10-13): xy, xy-2cell

\usepackage[all]{xy}
\xyoption{2cell}
\UseAllTwocells

\newenvironment{reference}{\comment}{\endcomment}
\newenvironment{slogan}{\comment}{\endcomment}
\newenvironment{history}{\comment}{\endcomment}
\theoremstyle{plain}
\newtheorem{theorem}[subsection]{Theorem}
\newtheorem{proposition}[subsection]{Proposition}
\newtheorem{lemma}[subsection]{Lemma}
\theoremstyle{definition}
\newtheorem{definition}[subsection]{Definition}
\newtheorem{example}[subsection]{Example}
\newtheorem{exercise}[subsection]{Exercise}
\newtheorem{situation}[subsection]{Situation}
\theoremstyle{remark}
\newtheorem{remark}[subsection]{Remark}
\newtheorem{remarks}[subsection]{Remarks}
\numberwithin{equation}{subsection}
\def\lim{\mathop{\mathrm{lim}}\nolimits}
\def\colim{\mathop{\mathrm{colim}}\nolimits}
\def\Spec{\mathop{\mathrm{Spec}}}
\def\Hom{\mathop{\mathrm{Hom}}\nolimits}
\def\Ext{\mathop{\mathrm{Ext}}\nolimits}
\def\SheafHom{\mathop{\mathcal{H}\!\mathit{om}}\nolimits}
\def\SheafExt{\mathop{\mathcal{E}\!\mathit{xt}}\nolimits}
\def\Sch{\mathit{Sch}}
\def\Mor{\mathop{\mathrm{Mor}}\nolimits}
\def\Ob{\mathop{\mathrm{Ob}}\nolimits}
\def\Sh{\mathop{\mathit{Sh}}\nolimits}
\def\NL{\mathop{N\!L}\nolimits}
\def\CH{\mathop{\mathrm{CH}}\nolimits}
\def\proetale{{pro\text{-}\acute{e}tale}}
\def\etale{{\acute{e}tale}}
\def\QCoh{\mathit{QCoh}}
\def\Ker{\mathop{\mathrm{Ker}}}
\def\Im{\mathop{\mathrm{Im}}}
\def\Coker{\mathop{\mathrm{Coker}}}
\def\Coim{\mathop{\mathrm{Coim}}}
\DeclareMathSymbol{\boxtimes}{\mathbin}{AMSa}{"02}
\def\QCohstack{\mathcal{QC}\!\mathit{oh}}
\def\Cohstack{\mathcal{C}\!\mathit{oh}}
\def\Spacesstack{\mathcal{S}\!\mathit{paces}}
\def\Quotfunctor{\mathrm{Quot}}
\def\Hilbfunctor{\mathrm{Hilb}}
\def\Curvesstack{\mathcal{C}\!\mathit{urves}}
\def\Polarizedstack{\mathcal{P}\!\mathit{olarized}}
\def\Complexesstack{\mathcal{C}\!\mathit{omplexes}}
\def\Pic{\mathop{\mathrm{Pic}}\nolimits}
\def\Picardstack{\mathcal{P}\!\mathit{ic}}
\def\Picardfunctor{\mathrm{Pic}}
\def\Deformationcategory{\mathcal{D}\!\mathit{ef}}
\usepackage{cleveref}
\crefname{originallemma}{original lemma}{original lemmas}
\Crefname{originallemma}{Original lemma}{Original lemmas}
\makeatletter
\newlabel{coherent-lemma-flat-base-change-cohomology}{{02KH}{source}{}{sourcepointer-1}{}}
\newlabel{coherent-lemma-flat-base-change-cohomology@cref}{{[originallemma][0][]02KH}{[1][0][]source}{}{}{}}
\makeatother
\crefname{originallemma}{original lemma}{original lemmas}
\Crefname{originallemma}{Original lemma}{Original lemmas}
\makeatletter
\newlabel{examples-defos-lemma-schemes-morphisms-smooth-to-base}{{0E3X}{source}{}{sourcepointer-2}{}}
\newlabel{examples-defos-lemma-schemes-morphisms-smooth-to-base@cref}{{[originallemma][0][]0E3X}{[1][0][]source}{}{}{}}
\makeatother
\crefname{originallemma}{original lemma}{original lemmas}
\Crefname{originallemma}{Original lemma}{Original lemmas}
\makeatletter
\newlabel{lemma-curves-qs-lfp}{{0DSS}{source}{}{sourcepointer-3}{}}
\newlabel{lemma-curves-qs-lfp@cref}{{[originallemma][0][]0DSS}{[1][0][]source}{}{}{}}
\makeatother
\crefname{originallemma}{original lemma}{original lemmas}
\Crefname{originallemma}{Original lemma}{Original lemmas}
\makeatletter
\newlabel{limits-lemma-descend-finite-presentation}{{01ZM}{source}{}{sourcepointer-4}{}}
\newlabel{limits-lemma-descend-finite-presentation@cref}{{[originallemma][0][]01ZM}{[1][0][]source}{}{}{}}
\makeatother
\crefname{originallemma}{original lemma}{original lemmas}
\Crefname{originallemma}{Original lemma}{Original lemmas}
\makeatletter
\newlabel{limits-lemma-eventually-proper}{{081F}{source}{}{sourcepointer-5}{}}
\newlabel{limits-lemma-eventually-proper@cref}{{[originallemma][0][]081F}{[1][0][]source}{}{}{}}
\makeatother
\crefname{originallemma}{original lemma}{original lemmas}
\Crefname{originallemma}{Original lemma}{Original lemmas}
\makeatletter
\newlabel{more-algebra-lemma-henselization-G-ring}{{07QR}{source}{}{sourcepointer-6}{}}
\newlabel{more-algebra-lemma-henselization-G-ring@cref}{{[originallemma][0][]07QR}{[1][0][]source}{}{}{}}
\makeatother
\crefname{originallemma}{original lemma}{original lemmas}
\Crefname{originallemma}{Original lemma}{Original lemmas}
\makeatletter
\newlabel{more-algebra-proposition-ubiquity-G-ring}{{07PX}{source}{}{sourcepointer-7}{}}
\newlabel{more-algebra-proposition-ubiquity-G-ring@cref}{{[originallemma][0][]07PX}{[1][0][]source}{}{}{}}
\makeatother
\crefname{originallemma}{original lemma}{original lemmas}
\Crefname{originallemma}{Original lemma}{Original lemmas}
\makeatletter
\newlabel{more-morphisms-lemma-describe-henselization}{{05KS}{source}{}{sourcepointer-8}{}}
\newlabel{more-morphisms-lemma-describe-henselization@cref}{{[originallemma][0][]05KS}{[1][0][]source}{}{}{}}
\makeatother
\crefname{originallemma}{original lemma}{original lemmas}
\Crefname{originallemma}{Original lemma}{Original lemmas}
\makeatletter
\newlabel{quot-lemma-curves-existence}{{0D58}{source}{}{sourcepointer-9}{}}
\newlabel{quot-lemma-curves-existence@cref}{{[originallemma][0][]0D58}{[1][0][]source}{}{}{}}
\makeatother
\crefname{originallemma}{original lemma}{original lemmas}
\Crefname{originallemma}{Original lemma}{Original lemmas}
\makeatletter
\newlabel{smoothing-theorem-approximation-property}{{07QY}{source}{}{sourcepointer-10}{}}
\newlabel{smoothing-theorem-approximation-property@cref}{{[originallemma][0][]07QY}{[1][0][]source}{}{}{}}
\makeatother
\crefname{originallemma}{original lemma}{original lemmas}
\Crefname{originallemma}{Original lemma}{Original lemmas}
\makeatletter
\newlabel{spaces-limits-lemma-descend-finite-presentation}{{07SK}{source}{}{sourcepointer-11}{}}
\newlabel{spaces-limits-lemma-descend-finite-presentation@cref}{{[originallemma][0][]07SK}{[1][0][]source}{}{}{}}
\makeatother
\crefname{originallemma}{original lemma}{original lemmas}
\Crefname{originallemma}{Original lemma}{Original lemmas}
\makeatletter
\newlabel{spaces-more-morphisms-lemma-algebraize-morphism}{{0A4Z}{source}{}{sourcepointer-12}{}}
\newlabel{spaces-more-morphisms-lemma-algebraize-morphism@cref}{{[originallemma][0][]0A4Z}{[1][0][]source}{}{}{}}
\makeatother
\crefname{originallemma}{original lemma}{original lemmas}
\Crefname{originallemma}{Original lemma}{Original lemmas}
\makeatletter
\newlabel{spaces-over-fields-lemma-codim-1-point-in-schematic-locus}{{0ADD}{source}{}{sourcepointer-13}{}}
\newlabel{spaces-over-fields-lemma-codim-1-point-in-schematic-locus@cref}{{[originallemma][0][]0ADD}{[1][0][]source}{}{}{}}
\makeatother
\crefname{originallemma}{original lemma}{original lemmas}
\Crefname{originallemma}{Original lemma}{Original lemmas}
\makeatletter
\newlabel{stacks-limits-lemma-representable-by-spaces-limit-preserving}{{0CMV}{source}{}{sourcepointer-14}{}}
\newlabel{stacks-limits-lemma-representable-by-spaces-limit-preserving@cref}{{[originallemma][0][]0CMV}{[1][0][]source}{}{}{}}
\makeatother
\title{Stacks, Moduli of Curves: lemma-contract-basic}
\author{The Stacks Project Authors}
\date{stacks-project:moduli-curves.tex:local-snapshot-20261012; GNU FDL 1.2 or later}
\begin{document}
\maketitle
\noindent\textit{Editorial note.} Statement span is lines 2986-2995 of sources/stacks-moduli-curves/pinned-original-source.tex; proof span is lines 2997-3124, the complete author proof (126 body lines).\par\medskip
\noindent\textit{Editorial note.} Navigation stubs: 14 cited-but-not-included labels are supplied as non-load-bearing pointers with their pinned tags where the snapshot tag index has them.\par\medskip
\noindent\textit{Editorial note.} extra\_packages is empty on purpose: the retained preamble already loads xy with options and a second declaration clashes in the assembler wrapper.\par\medskip
\noindent\textit{Editorial note.} Text is retained verbatim, including the author's own typos and placeholders.\par\medskip
\noindent\textit{Editorial note.} Queue index 0 of queue\_moduli-curves.json (pb=5310, stmt 2986-2995, proof 2997-3124).\par\medskip
\noindent\textit{Allowed standard prerequisites (editorial).}\begin{itemize}
\item Explicit prerequisite boundary. The cited inputs are prerequisites; the argument in the retained span is not replaced by any cited result.
\end{itemize}\medskip
\ifdefined\hypertarget\hypertarget{sourcepointer-1}{}\fi
\noindent\textit{Source reference.} Original lemma 02KH: sources/stacks-moduli-curves/pinned-original-source.tex, source lines 3124--3124. Cited by the retained statement or proof but not part of the collected spans; supplied only as a navigation pointer. The argument does not depend on its content.\par\medskip
\ifdefined\hypertarget\hypertarget{sourcepointer-2}{}\fi
\noindent\textit{Source reference.} Original lemma 0E3X: sources/stacks-moduli-curves/pinned-original-source.tex, source lines 3124--3124. Cited by the retained statement or proof but not part of the collected spans; supplied only as a navigation pointer. The argument does not depend on its content.\par\medskip
\ifdefined\hypertarget\hypertarget{sourcepointer-3}{}\fi
\noindent\textit{Source reference.} Original lemma 0DSS: sources/stacks-moduli-curves/pinned-original-source.tex, source lines 655--658. Cited by the retained statement or proof but not part of the collected spans; supplied only as a navigation pointer. The argument does not depend on its content.\par\medskip
\ifdefined\hypertarget\hypertarget{sourcepointer-4}{}\fi
\noindent\textit{Source reference.} Original lemma 01ZM: sources/stacks-moduli-curves/pinned-original-source.tex, source lines 3124--3124. Cited by the retained statement or proof but not part of the collected spans; supplied only as a navigation pointer. The argument does not depend on its content.\par\medskip
\ifdefined\hypertarget\hypertarget{sourcepointer-5}{}\fi
\noindent\textit{Source reference.} Original lemma 081F: sources/stacks-moduli-curves/pinned-original-source.tex, source lines 3124--3124. Cited by the retained statement or proof but not part of the collected spans; supplied only as a navigation pointer. The argument does not depend on its content.\par\medskip
\ifdefined\hypertarget\hypertarget{sourcepointer-6}{}\fi
\noindent\textit{Source reference.} Original lemma 07QR: sources/stacks-moduli-curves/pinned-original-source.tex, source lines 3124--3124. Cited by the retained statement or proof but not part of the collected spans; supplied only as a navigation pointer. The argument does not depend on its content.\par\medskip
\ifdefined\hypertarget\hypertarget{sourcepointer-7}{}\fi
\noindent\textit{Source reference.} Original lemma 07PX: sources/stacks-moduli-curves/pinned-original-source.tex, source lines 3124--3124. Cited by the retained statement or proof but not part of the collected spans; supplied only as a navigation pointer. The argument does not depend on its content.\par\medskip
\ifdefined\hypertarget\hypertarget{sourcepointer-8}{}\fi
\noindent\textit{Source reference.} Original lemma 05KS: sources/stacks-moduli-curves/pinned-original-source.tex, source lines 3124--3124. Cited by the retained statement or proof but not part of the collected spans; supplied only as a navigation pointer. The argument does not depend on its content.\par\medskip
\ifdefined\hypertarget\hypertarget{sourcepointer-9}{}\fi
\noindent\textit{Source reference.} Original lemma 0D58: sources/stacks-moduli-curves/pinned-original-source.tex, source lines 3124--3124. Cited by the retained statement or proof but not part of the collected spans; supplied only as a navigation pointer. The argument does not depend on its content.\par\medskip
\ifdefined\hypertarget\hypertarget{sourcepointer-10}{}\fi
\noindent\textit{Source reference.} Original lemma 07QY: sources/stacks-moduli-curves/pinned-original-source.tex, source lines 3124--3124. Cited by the retained statement or proof but not part of the collected spans; supplied only as a navigation pointer. The argument does not depend on its content.\par\medskip
\ifdefined\hypertarget\hypertarget{sourcepointer-11}{}\fi
\noindent\textit{Source reference.} Original lemma 07SK: sources/stacks-moduli-curves/pinned-original-source.tex, source lines 3124--3124. Cited by the retained statement or proof but not part of the collected spans; supplied only as a navigation pointer. The argument does not depend on its content.\par\medskip
\ifdefined\hypertarget\hypertarget{sourcepointer-12}{}\fi
\noindent\textit{Source reference.} Original lemma 0A4Z: sources/stacks-moduli-curves/pinned-original-source.tex, source lines 3124--3124. Cited by the retained statement or proof but not part of the collected spans; supplied only as a navigation pointer. The argument does not depend on its content.\par\medskip
\ifdefined\hypertarget\hypertarget{sourcepointer-13}{}\fi
\noindent\textit{Source reference.} Original lemma 0ADD: sources/stacks-moduli-curves/pinned-original-source.tex, source lines 3124--3124. Cited by the retained statement or proof but not part of the collected spans; supplied only as a navigation pointer. The argument does not depend on its content.\par\medskip
\ifdefined\hypertarget\hypertarget{sourcepointer-14}{}\fi
\noindent\textit{Source reference.} Original lemma 0CMV: sources/stacks-moduli-curves/pinned-original-source.tex, source lines 3124--3124. Cited by the retained statement or proof but not part of the collected spans; supplied only as a navigation pointer. The argument does not depend on its content.\par\medskip
\begin{lemma}
\label{lemma-contract-basic}
Let $f : X \to S$ be a family of curves. Let $s \in S$ be a point.
Let $h_0 : X_s \to Y_0$ be a morphism to a proper scheme $Y_0$ over $\kappa(s)$
such that $h_{0, *}\mathcal{O}_{X_s} = \mathcal{O}_{Y_0}$ and
$R^1h_{0, *}\mathcal{O}_{X_s} = 0$. Then there exist an elementary
\'etale neighbourhood $(U, u) \to (S, s)$, a family of curves $Y \to U$,
and a morphism $h : X_U \to Y$ over $U$ whose fibre in $u$
is isomorphic to $h_0$.
\end{lemma}

\begin{proof}
We first do some reductions; we urge the reader to skip ahead.
The question is local on $S$, hence we may assume $S$ is affine.
Write $S = \lim S_i$ as a cofiltered limit of affine schemes $S_i$
of finite type over $\mathbf{Z}$.
For some $i$ we can find a family of curves $X_i \to S_i$
whose base change is $X \to S$. This follows from
Lemma \ref{lemma-curves-qs-lfp} and
Limits of Stacks, Lemma
\ref{stacks-limits-lemma-representable-by-spaces-limit-preserving}.
Let $s_i \in S_i$ be the image of $s$. Observe that
$\kappa(s) = \colim \kappa(s_i)$ and that $X_s$ is a scheme
(Spaces over Fields, Lemma
\ref{spaces-over-fields-lemma-codim-1-point-in-schematic-locus}).
After increasing $i$ we may assume there exists a morphism
$h_{i, 0} : X_{i, s_i} \to Y_i$
of finite type schemes over $\kappa(s_i)$ whose base change to
$\kappa(s)$ is $h_0$, see
Limits, Lemma \ref{limits-lemma-descend-finite-presentation}.
After increasing $i$ we may assume $Y_i$ is proper over $\kappa(s_i)$, see
Limits, Lemma \ref{limits-lemma-eventually-proper}.
Let $g_{i, 0} : Y_0 \to Y_{i, 0}$ be the projection. Observe that
this is a faithfully flat morphism as the base change of
$\Spec(\kappa(s)) \to \Spec(\kappa(s_i))$.
By flat base change we have
$$
h_{0, *}\mathcal{O}_{X_s} = g_{i, 0}^*h_{i, 0, *}\mathcal{O}_{X_{i, s_i}}
\quad\text{and}\quad
R^1h_{0, *}\mathcal{O}_{X_s} = g_{i, 0}^*Rh_{i, 0, *}\mathcal{O}_{X_{i, s_i}}
$$
see Cohomology of Schemes, Lemma
\ref{coherent-lemma-flat-base-change-cohomology}.
By faithful flatness we see that $X_i \to S_i$, $s_i \in S_i$, and
$X_{i, s_i} \to Y_i$ satisfies all the assumptions of the lemma.
This reduces us to the case discussed in the next paragraph.

\medskip\noindent
Assume $S$ is affine of finite type over $\mathbf{Z}$.
Let $\mathcal{O}_{S, s}^h$ be the henselization of the local
ring of $S$ at $s$. Observe that $\mathcal{O}_{S, s}^h$ is a G-ring by
More on Algebra, Lemma \ref{more-algebra-lemma-henselization-G-ring} and
Proposition \ref{more-algebra-proposition-ubiquity-G-ring}.
Suppose we can construct a family of curves
$Y' \to \Spec(\mathcal{O}_{S, s}^h)$ and a morphism
$$
h' : X \times_S \Spec(\mathcal{O}_{S, s}^h) \longrightarrow Y'
$$
over $\Spec(\mathcal{O}_{S, s}^h)$ whose base change to the closed
point is $h_0$. This will be enough. Namely, first we use that
$$
\mathcal{O}_{S, s}^h = \colim_{(U, u)} \mathcal{O}_U(U)
$$
where the colimit is over the filtered category of 
elementary \'etale neighbourhoods (More on Morphisms, Lemma
\ref{more-morphisms-lemma-describe-henselization}).
Next, we use again that given $Y'$ we can descend it to
$Y \to U$ for some $U$ (see references given above).
Then we use
Limits, Lemma \ref{limits-lemma-descend-finite-presentation}
to descend $h'$ to some $h$. This reduces us to the case
discussed in the next paragraph.

\medskip\noindent
Assume $S = \Spec(\Lambda)$ where $(\Lambda, \mathfrak m, \kappa)$
is a henselian Noetherian local G-ring and $s$ is the closed point of $S$.
Recall that the map
$$
\Deformationcategory_{X_s \to Y_0} \to \Deformationcategory_{X_s}
$$
is an equivalence, see Deformation Problems, Lemma
\ref{examples-defos-lemma-schemes-morphisms-smooth-to-base}.
(This is the only important step in the proof; everything else
is technique.) Denote $\Lambda^\wedge$ the $\mathfrak m$-adic completion.
The pullbacks $X_n$ of $X$ to $\Lambda/\mathfrak m^{n + 1}$ define
a formal object $\xi$ of $\Deformationcategory_{X_s}$ over $\Lambda^\wedge$.
From the equivalence we obtain a formal object
$\xi'$ of $\Deformationcategory_{X_s \to Y_0}$ over $\Lambda^\wedge$.
Thus we obtain a huge commutative diagram
$$
\xymatrix{
\ldots \ar[r] &
X_n \ar[r] \ar[d] &
X_{n - 1} \ar[r] \ar[d] &
\ldots \ar[r] &
X_s \ar[d] \\
\ldots \ar[r] &
Y_n \ar[r] \ar[d] &
Y_{n - 1} \ar[r] \ar[d] &
\ldots \ar[r] &
Y_0 \ar[d] \\
\ldots \ar[r] &
\Spec(\Lambda/\mathfrak m^{n + 1}) \ar[r] &
\Spec(\Lambda/\mathfrak m^n) \ar[r] &
\ldots \ar[r] &
\Spec(\kappa)
}
$$
The formal object $(Y_n)$ comes from a family of curves
$Y' \to \Spec(\Lambda^\wedge)$ by
Quot, Lemma \ref{quot-lemma-curves-existence}.
By More on Morphisms of Spaces, Lemma
\ref{spaces-more-morphisms-lemma-algebraize-morphism}
we get a morphism $h' : X_{\Lambda^\wedge} \to Y'$ inducing
the given morphisms $X_n \to Y_n$ for all $n$
and in particular the given morphism $X_s \to Y_0$.

\medskip\noindent
To finish we do a standard algebraization/approximation argument.
First, we observe that we can find a finitely generated $\Lambda$-subalgebra
$\Lambda \subset A \subset \Lambda^\wedge$, a family of curves
$Y'' \to \Spec(A)$ and a morphism $h'' : X_A \to Y''$ over $A$
whose base change to $\Lambda^\wedge$ is $h'$.
This is true because $\Lambda^\wedge$ is the filtered colimit
of these rings $A$ and we can argue as before using
that $\Curvesstack$ is locally of finite presentation
(which gives us $Y''$ over $A$ by
Limits of Stacks, Lemma
\ref{stacks-limits-lemma-representable-by-spaces-limit-preserving})
and using 
Limits of Spaces, Lemma \ref{spaces-limits-lemma-descend-finite-presentation}
to descend $h'$ to some $h''$.
Then we can apply the approximation property for G-rings
(in the form of Smoothing Ring Maps, Theorem
\ref{smoothing-theorem-approximation-property})
to find a map $A \to \Lambda$ which induces the same map
$A \to \kappa$ as we obtain from $A \to \Lambda^\wedge$.
Base changing $h''$ to $\Lambda$ the proof is complete.
\end{proof}

\end{document}
